
\section{Exactness and other boundaries}\label{sec:exactness}

Theorem~\ref{thm:boundary} concerns rational configurations and sequence
existence at every positive tolerance. Its word-against-pair hardness
reduction uses growing probe sets. In the single-word class, signed cores
classify the real components of the budgeted eigenvalue locus
(Corollary~\ref{cor:real}), and Proposition~\ref{prop:separation} gives an
explicit separation cutoff for rational eigenvalue vectors with different
signed cores. With one word-against-pair test, real and rational verdicts
can differ. Without probes, only part of that class is decided, whereas
many such tests give PSPACE-hardness. Appendix~\ref{app:exactness-proofs}
has the proofs.

% ---------------------------------------------------------------------------
\subsection{Real versus rational reachability}\label{sec:real-rational}

The real--rational agreement of Corollary~\ref{cor:real} can fail with one
word-against-pair test, even in the following models. They have one register value, identity reversal, and one full-support profile at level~$0$ holding every listed test, whose word tests form the only admissible set.

\begin{proposition}[One probe]\label{prop:one-probe}
Take a binary base with one probe, $\kappa=1/2$, $\delta=1/9$, $H=2$ and
labels $\mathsf u,\mathsf h,\mathsf x,\mathsf y,\mathsf x_2,\mathsf
y_2,\mathsf x_4,\mathsf y_4$. Hold the calibrations $\alpha_{\mathsf
u}=1/8$ and $\alpha_{\mathsf h}=1/4$ and the word tests
\[
  \mathsf u\sim\{\mathsf x_4,\mathsf y_4\},\quad
  \mathsf x_2\mathsf h\sim\mathsf{xx},\quad
  \mathsf x_4\mathsf h\sim\mathsf x_2\mathsf x_2,\quad
  \mathsf y_2\mathsf h\sim\mathsf{yy},\quad
  \mathsf y_4\mathsf h\sim\mathsf y_2\mathsf y_2 .
\]
With $x,y$ the eigenvalues of $\mathsf x,\mathsf y$, there are unique
feasible configurations $A$ and $B$ with $(x,y)=(1/4,0)$ and $(0,1/4)$. A
continuous path in the real feasible set joins them, but no
reconfiguration sequence at tolerance $1/72$ does.
\end{proposition}

These tests force $x^4+y^4=4^{-4}$, whose rational points lie only on the
axes by Fermat's theorem for exponent four.

\begin{proposition}[No probe and no calibration]\label{prop:zero-probe-fermat}
Take one base state, $\kappa=0$, $\delta=1/40$, $H=4$, labels
$\mathsf u,\mathsf h,\mathsf x,\mathsf y,\mathsf x_4,\mathsf y_4$, the
word tests $\mathsf u\sim\{\mathsf x_4,\mathsf y_4\}$,
$\mathsf x_4\mathsf{hhh}\sim\mathsf{xxxx}$ and
$\mathsf y_4\mathsf{hhh}\sim\mathsf{yyyy}$, and the probability-ratio test
$\theta(\mathsf u)=17\,\theta(\mathsf h)$. With
$x=\theta_{\mathsf x}/\theta_{\mathsf h}$ and
$y=\theta_{\mathsf y}/\theta_{\mathsf h}$, there are unique feasible
configurations $A$ and $B$ with $(x,y)=(1,2)$ and $(2,1)$. A real path
joins them, but $B$ is not finely
reachable from~$A$.
\end{proposition}

Here $x^4+y^4=17$ has finitely many rational points, by Darmon and
Granville~\cite[Theorem~2]{darmongranville}. So testing component membership in the full real feasible set, for
instance by a roadmap query~\cite{canny93,basupollackroy2006book}, can give
the wrong rational verdict. This mismatch does not exclude roadmaps within
algorithms with further arithmetic steps.
By Proposition~\ref{prop:zero-probe-fermat}, the no-calibration
reachability conclusion of Theorem~\ref{thm:no-calibration} does not
extend to models with a word-against-pair test.

% ---------------------------------------------------------------------------
\subsection{Separation}\label{sec:separation}

\begin{proposition}[Separation cutoff]\label{prop:separation}
\begin{enumerate}
\item[(a)] Fix a budgeted signed-binomial system on $J$ with $n=|J|\ge1$
and $\FQ$ bounded. Let $\mathcal L$ be the binary length of its coordinates, tests
and budgets (not endpoints), and $D_0$ as in Theorem~\ref{thm:certificates}.
Let $D_1$ be the maximum of $1$ and the binary lengths of the numerator
magnitudes and denominators of all test coefficients and budget bounds.
For
\[
  D=1+D_1\,n\,n!\,D_0^{\,n}+D_1\,n^3(n!)^2D_0^{\,2n+1},
\]
feasible rational points with different signed cores are at sup distance
at least $2^{-D}$. Here $\log_2D$ is polynomial in~$\mathcal L$, while $2^{-D}$ has
a denominator of $D+1$ binary digits.
\item[(b)] Let $A$ and $B$ be feasible rational configurations of a
single-word model with $\SC_{p_A}(\alpha_A)\neq\SC_{p_A}(\alpha_B)$, with
$D$ computed for $\Phi_{p_A}$ and $N=|X||R|$. Then no reconfiguration
sequence at tolerance $\delta2^{-D}/(2N+1)$ leads from $A$ to~$B$.
\item[(c)] For $s\ge1$, the system $x_0=1/2$, $x_{i+1}=x_i^2$
$(0\le i<s)$, $z^2=x_s^2$ with budgets $|x_i|\le1$ and $|z|\le1$ has
exactly two feasible rational points, with different signed cores, at
sup distance $2^{1-2^s}$. A one-probe single-word model realizes them
at one profile, $2^{-2^s-1}/(s+3)$ apart in its tables.
\end{enumerate}
\end{proposition}

For a single-word model with at least one probe, Theorem~\ref{thm:directed}
and part~(b) give equivalence between fine reachability and sequence
existence at $\delta2^{-D}/(2N+1)$, with $D$ computed at the source
profile. If $p_A\npreceq p_B$, no sequence exists at any tolerance; with
no probes, the directed profile order alone decides fine reachability. By~(c) and its binary-exponent variant
(Appendix~\ref{app:separation}), valid cutoffs can need exponentially many
binary digits. Jeronimo,
Perrucci and Tsigaridas bound the distance between disjoint components of
basic closed integer semialgebraic sets, one of them compact, at this
doubly exponential scale~\cite[Theorem~1.2]{jeronimoperruccitsigaridas}.
Their Example~3~\cite{jeronimoperruccitsigaridas} with both parameters
two is the system of~(c) up to indexing and redundant budgets.
Our contribution is the direct bound of~(a) for signed-core classes of
rational points under rational boundedness, and its Markov transfer~(b).

% ---------------------------------------------------------------------------
\subsection{Zero probes}\label{sec:zero-probes}

Section~\ref{sec:quartic} leaves zero probes open. Even with no probes,
labels $\mathsf a,\mathsf b,\mathsf c$, probability floor $1/8$,
$\mathsf c\sim\{\mathsf a,\mathsf b\}$ and
$\theta(\mathsf{ab})=\tfrac3{16}\theta(\mathsf{cc})$ allow only
$(\theta_{\mathsf a},\theta_{\mathsf b})=(1/8,3/8)$ and $(3/8,1/8)$, a
nonconvex set. For exactly one word-against-pair test
$w_0\sim\{w_1,w_2\}$, if $p_A$ does not hold it, $B$ is finely reachable
from~$A$ exactly when $p_A\preceq p_B$.
Let $t=\theta(w_1)/\theta(w_0)$. In the notation of
Proposition~\ref{prop:zero-probes} (Appendix~\ref{app:exactness-proofs}),
the image rank is $\dim_{\RR}(f,g)(\ker_{\RR}\Psi)$. In rank two, the
case is hyperbolic when $1/d_0+1/d_1+1/d_\infty<1$. If the source holds the test,
image rank at most one and hyperbolic rank two each give a polynomial-time
criterion: $B$ is finely reachable from $A$ exactly when $p_A\preceq p_B$
and $t_A=t_B$ (Proposition~\ref{prop:zero-probes}). In the remaining
rank-two case, $B$ is finely reachable from $A$ exactly when
$p_A\preceq p_B$ and the closure of the feasible rational ratios at $p_A$
contains the closed interval between $t_A$ and~$t_B$
(Proposition~\ref{prop:zero-probe-closure}). Deciding this containment in
the remaining case is open. The directed order matters even when both
ratios agree. For a rank-one example, take one base state, one register value, label
probabilities $(1/2,1/4,1/4)$ on $(\mathsf L,\mathsf A,\mathsf B)$ and
probability floor $1/4$. Both configurations hold
$\mathsf L\sim\{\mathsf A,\mathsf B\}$ and $\mathsf A\sim\mathsf B$, and
only the source holds $\theta_{\mathsf A}/\theta_{\mathsf B}=1$. With all
other profile data equal, the profiles are comparable and both ratios $t$
are $1/2$, but no directed sequence can drop the source's held
probability-ratio test.
The paper establishes neither hardness with zero probes and exactly one
word-against-pair test nor a general decision bound for fixed probe
counts.

% ---------------------------------------------------------------------------
\subsection{The one-state subclass}\label{sec:one-state}

\begin{proposition}[The one-state subclass]\label{prop:one-state}
Consider fixed-profile Fine Reachability with one base state and one
register value, hence no probes. Take $\kappa=0$ and identity label
reversal (the \defterm{one-state subclass}).
\begin{enumerate}
\item[(a)] With words of length two, a positive probability floor and a
growing number of word-against-pair tests, it is PSPACE-hard under
polynomial-time many-one reductions with distinct endpoints. %
The same reduction proves PSPACE-hardness at the supplied tolerance
$2\delta$. On each output, sequence existence has the same verdict at
every positive tolerance at most~$2\delta$.
\item[(b)] If all held equations are linear in the label probabilities,
every feasible pair is finely reachable. This includes the case in which
all held word tests and probability-ratio tests have length one. The same
conclusion holds if every held word test is word-against-word, by
homogeneity and positivity of the ratio constants.
\item[(c)] Rational-zero existence for sparsely binary-encoded integer
polynomials reduces in polynomial time to the problem of~(a), also at its
supplied tolerances; on each output, the real feasible set at the profile
is path-connected.
\end{enumerate}
\end{proposition}

Part~(a) encodes 3-CNF two-solution
connectivity~\cite[Lemma~3.6]{gopalan} through Schnider and Weber's
cubical complex~\cite[Definition~4]{schniderweber}, and does not give
Theorem~\ref{thm:boundary}(ii).
Word tests of length one alone give linear rate equations in this
subclass, as do probability-ratio tests of length one. Some length-two
probability ratios also reduce to linear equations: a held ratio
$\theta(\mathsf{ab})=c\,\theta(\mathsf{cb})$ reduces to
$\theta_{\mathsf a}=c\,\theta_{\mathsf c}$ because
$\theta_{\mathsf b}>0$.
A length-one word-against-pair test together with a length-two
probability-ratio test can give a no-instance, as in the nonconvex example
of Section~\ref{sec:zero-probes}.

Part~(c) compiles the switch of Section~\ref{sec:rational-roots}, so a
PSPACE upper bound for~(a) would put rational-zero existence in PSPACE.
Chonev's Remark~7~\cite{chonev} encodes polynomial threshold comparisons by
reachability probabilities. Winkler et al., Lemma~20
\cite{winklerjungesperezkatoen}, give pointwise equivalences on the
parameter box for five comparison relations, including equality. In the
extended article, Junges, Katoen, Pérez and Winkler's
Proposition~11~\cite{jungeskatoenperezwinkler} realizes an affine
rescaling of a polynomial as a reachability probability.
